\section{Open questions}
\label{sec:open}

Five open questions surfaced by this replication (see also
\texttt{report/open\_questions.json} for the machine-readable version):

\begin{enumerate}

\item \textbf{Non-CSS extension.} Does the inplace Y-basis construction extend
to non-CSS codes (color codes, XZZX, Bacon--Shor), or is diagonal twist-defect
fusion specifically tied to the rotated-surface CSS structure? The paper's
construction relies on twist defects that live on the standard rotated surface
code's XZ boundary structure; whether an analogous ``inplace'' non-braid
Y-basis access exists for non-CSS or twisted-CSS codes is not addressed.
\emph{Next steps:} attempt an analogous construction on the color code
(Steane hex lattice), which supports transversal H and therefore has a
natural Y-basis question; test on XZZX where noise-biased asymmetry may allow
an even shorter inplace path.

\item \textbf{Lattice-surgery composition for native $Y_1Y_2$ parity.} Can the
inplace-Y measurement be composed with lattice surgery for direct multi-patch
Y-basis parity checks without going through XZ-decomposition? Lattice surgery
is standardly defined for X-type and Z-type joins. A native Y-type join built
on top of inplace-Y would eliminate one layer of XZ-decomposition overhead in
factory circuits, potentially compounding per-Y savings. \emph{Next steps:}
design a merge protocol between two inplace-Y patches, simulate the merged
$Y_1Y_2$ parity extraction in \texttt{stim}, and compare LER vs the
XZ-decomposed baseline (measure $X_1X_2$ then $Z_1Z_2$, combine).

\item \textbf{End-to-end magic-state distillation gains.} What is the actual
factory-level saving when inplace-Y is used inside a magic-state distillation
protocol (e.g.\ 15-to-1 or 116-to-12) versus the standard Y-via-distillation
baseline? The paper motivates its construction by noting factories consume
hundreds of Y-basis measurements per T state, but does not simulate a full
factory. \emph{Next steps:} compose the released Y-basis circuit into a
15-to-1 factory (Litinski's compact layout), sample end-to-end, compare
T-state error rate and qubit-round volume against the standard baseline at
matched physical $p$.

\item \textbf{Biased-noise scaling.} How does the inplace-Y LER scale under
biased noise ($Z$-biased with bias $\eta = 100$ or $1000$) rather than
SI1000-depolarizing? The inplace construction concentrates twist defects along
a diagonal --- a locally denser region --- and biased-noise effective distance
can differ sharply from depolarizing distance. \emph{Next steps:} rebuild the
released \texttt{.stim} circuits with a biased-depolarizing channel at the
same locations, sample LER at $\eta \in \{1, 10, 100, 1000\}$ for
$d \in \{5, 7, 9\}$, check whether the inplace-vs-braid relative advantage is
preserved or inverted.

\item \textbf{Hardware-native compilation.} Which superconducting and
neutral-atom hardware platforms can execute the diagonal twist-defect fusion
natively, and what is the true physical-qubit and native-gate overhead once
platform connectivity and gate set are honored? Ideal square connectivity is
assumed; heavy-hex, linear-nearest-neighbor, and neutral-atom rearrangement
systems may compile poorly. \emph{Next steps:} take the released Y-basis
circuit, run naive SWAP-network compilation onto (i) heavy-hex,
(ii) linear-nearest-neighbor, (iii) neutral-atom rearrangeable connectivity,
measure depth and two-qubit-gate blow-up, and re-simulate LER with the
corresponding platform-specific noise model.

\end{enumerate}
